\documentclass[11pt]{article}

\usepackage[a4paper,margin=2.4cm]{geometry}
\usepackage{amsmath,amssymb}
\usepackage{booktabs}
\usepackage{siunitx}
\usepackage{xcolor}
\usepackage[hidelinks]{hyperref}
\usepackage{enumitem}

\newcommand{\bplus}{\beta^{+}}
\newcommand{\Cl}[1]{{}^{#1}\mathrm{C}}
\newcommand{\Ox}[1]{{}^{#1}\mathrm{O}}
\newcommand{\Nt}[1]{{}^{#1}\mathrm{N}}

\title{PET detector simulation and comparison (Stages B--C)}
\author{Working document}
\date{\today}

\begin{document}
\maketitle

% NOTE: cut verbatim from the original chain spec (ptcryspg4/docs/simulate_pt_pet.tex).
% Cross-references to the source stage (Eq. off, Sec. production / timedecay) point
% to the upstream document and will be reconnected when this is reworked.

\section{Stage B: detector --- from annihilation events to the coincidence list}
\label{sec:detector}

This stage runs \textbf{once per candidate PET}. It reads the \emph{same} sampled
annihilation events and produces a \textbf{coincidence list} -- the output of the
simulation and the only input to reconstruction (Sec.~\ref{sec:recon}). It is an
\textbf{analytic $\gamma$-transport Monte Carlo} (in Julia, a separate
repository), not Geant4: the simple ring geometry and the 511\,keV physics are
done directly, which is faster, transparent, and avoids a per-detector Geant4
build.
\begin{enumerate}[nosep]
\item \textbf{Detector simulation (analytic MC):} the two \SI{511}{keV} photons
propagate through the phantom (treated as water; attenuation and Compton scatter
by Klein--Nishina) and into the crystal. The interaction depth in the crystal
gives the depth-of-interaction; energy resolution $\sigma_E$, position/DOI
resolution and time resolution $\sigma_t$ are applied to the interaction; the
ring acceptance is a ray--cylinder intersection.
\item \textbf{Coincidence sorting:} energy window, coincidence time window;
classify true/scatter/random. In the in-room scenario the in-beam
prompt-$\gamma$/neutron background is absent; the relevant backgrounds are random
and patient-scattered coincidences, plus -- for a LYSO candidate -- the intrinsic
${}^{176}\mathrm{Lu}$ singles floor, injected as a volumetric source in the
crystals (Sec.~\ref{sec:compare}).
\item \textbf{Coincidence list (the output).} Each accepted coincidence is written
with everything a list-mode reconstructor needs:
\begin{center}
the two \textbf{DOI-resolved 3-D hit positions} (or crystal IDs + a known geometry),
the two \textbf{energies}, and the two \textbf{timestamps} (for TOF).
\end{center}
Truth flags (true/scatter/random) may be appended for our own analysis; the
reconstructor does not require them.
\end{enumerate}

\section{Reconstruction (independent stage)}
\label{sec:recon}

Reconstruction is deliberately \textbf{decoupled} from the simulation: it consumes
only the coincidence list, so the toolkit (a list-mode MLEM/OSEM of our own, or an
external package such as CASToR or STIR) can be chosen or replaced later without
touching stages~A or~B. For the present plan it suffices to fix the requirements:
\begin{itemize}[nosep]
\item \textbf{List-mode} (not sinogram-binned): the measurement is photon-starved
and time-resolved (per-frame isotope separation), and the events carry per-event
DOI.
\item \textbf{DOI in the system model:} oblique lines of response are exactly where
it prevents parallax blur, so the reconstructor must use the 3-D hit positions, not
a fixed crystal-face radius.
\item \textbf{TOF-aware} when the timing axis is being compared.
\end{itemize}
The output is the 3-D activity image and, in particular, the 1-D depth profile
along the beam axis. Reconstruction must always be performed (not merely a
projection-space comparison), since the resolution, scatter-rejection and parallax
differences between detectors only fully manifest in the image.

\section{Comparing the detectors}
\label{sec:compare}

\paragraph{Same source, different responses.}
Every candidate PET is fed the identical annihilation events. Any difference in
the reconstructed images therefore comes from the detector (geometry, sensitivity,
$\sigma_E$, $\sigma_t$, DOI), not from the source. This is what makes the Geant4
cross-section weakness irrelevant to the \emph{ranking}: a systematic error in
$P_j(\mathbf{r})$ is common to all detectors and cancels. The file can later be
regenerated from a better production map without re-running any detector
simulation.

\paragraph{Figure of merit: range precision.}
The clinically meaningful quantity is the position of the activity distal
fall-off, which tracks the proton range (it sits a few mm \emph{proximal} to the
Bragg peak, because production needs above-threshold protons). For each detector:
\begin{enumerate}[nosep]
\item fix the count budget from Eq.~\eqref{eq:off} at $t_{\mathrm{del}}\approx\SI{120}{s}$;
for a single \SI{1}{Gy} field this is the photon-starved regime of $10^{7}$--$10^{8}$
decays, i.e.\ only $\sim\!10^{5}$--$10^{6}$ collected coincidences after attenuation
and efficiency;
\item generate $M$ independent Poisson realisations at that budget;
\item reconstruct each, fit the distal fall-off, and report the spread
$\sigma(\text{range})$ in mm.
\end{enumerate}
The headline plot is $\sigma(\text{range})$ \textbf{vs.\ counts} (equivalently vs.\
dose, or vs.\ delay $t_{\mathrm{del}}$, which sets how much $\Ox{15}$ survives) per
detector. Secondary metrics: recovered axial/lateral resolution, scatter fraction
retained inside the energy window, and parallax blur at large axial angle.

\paragraph{What to keep fixed vs.\ vary.}
Fixed across detectors: the annihilation-event source, the phantom, the count
budget, and the reconstruction algorithm family. Varied: only the detector. The
discriminators that matter for in-room range verification of this photon-starved,
scatter-dominated, $\Ox{15}$-rich signal -- and which the matrix is built to
exercise -- are:
\begin{itemize}[nosep]
\item \textbf{Solid-angle sensitivity (AFOV).} The leading requirement. Sweep the
axial length $L$; the geometric acceptance for a centred source scales roughly as
$L/\sqrt{L^2+R^2}$, the dominant lever on how many of the scarce pairs are
collected.
\item \textbf{Energy resolution $\sigma_E$.} Sets how tightly the \SI{511}{keV}
window can be drawn, hence how much patient scatter is rejected on the
\emph{oblique} coincidences a long AFOV adds. Compare e.g.\ $\sim\!6\%$ (cryogenic
CsI) vs.\ $\sim\!10\%$ (LYSO) vs.\ poor (BGO).
\item \textbf{Depth-of-interaction.} Controls parallax on those same oblique LORs.
Compare DOI-capable monolithic crystals against fixed-depth pixels.
\item \textbf{Intrinsic radioactivity.} Inject the ${}^{176}\mathrm{Lu}$ singles
floor for the LYSO arm (none for CsI/BGO); at sub-MBq signal it raises randoms and
competes with the signal, so it is part of the comparison, not a detail.
\item \textbf{Timing / TOF.} Treated as a \emph{low-value} axis here: include a
TOF-on variant to demonstrate diminishing returns for a one-dimensional
distal-edge task, rather than as a discriminator the candidate relies on.
\end{itemize}
\textbf{Candidate set:} cryogenic monolithic CsI (the CRYSP baseline) against
conventional LYSO and BGO scanners, plus the CRYSP variants -- room-temperature
CsI(Tl), cryogenic BGO, and the BGO-core/CsI-wing hybrid -- each represented by its
$(\text{geometry},\,\sigma_E,\,\text{DOI resolution},\,\text{intrinsic background})$.
Since all see the identical source, the resulting $\sigma(\text{range})$ vs.\
counts curves isolate exactly these detector properties.

\section{Scope and honest caveats}
\label{sec:scope}

\begin{itemize}[nosep]
\item \textbf{Homogeneous phantom} omits tissue heterogeneity, hence no
$\Ox{15}/\Cl{11}$ spatial split and no washout (deliberate; this is a detector
study). A two-region phantom (e.g.\ a bone-like insert) can be added later to
probe density-interface artifacts.
\item \textbf{Geant4 cross sections} bias absolute counts and the production
shape; the relative comparison is robust, the absolute range precision is
indicative.
\item \textbf{In-room scope.} Only the in-room modality is treated (full-ring
scanner, $t_{\mathrm{del}}\gtrsim\SI{120}{s}$); in-beam and offline are out of
scope. At this short delay the source is $\Ox{15}$-dominated.
\item \textbf{Positron range} ($\Ox{15}$, $\Cl{10}$) sets a physical resolution
floor; a high-performance PET may saturate against it -- a result to report, not
to hide. It is captured by the PROD$\to$ANH displacement.
\item \textbf{Discarded prompt $\gamma$ / secondary neutrons} are small omissions,
acceptable for a first cut and listed here for honesty.
\end{itemize}

\end{document}
